% A LaTeX template for MSc Thesis submissions to
% Politecnico di Milano (PoliMi) - School of Industrial and Information Engineering
%
% Adapted by Salvatore Conte, 2026
% Original template: S. Bonetti, A. Gruttadauria, G. Mescolini, A. Zingaro
% e-mail: template-tesi-ingind@polimi.it

\documentclass{Configuration_Files/PoliMi3i_thesis}

%------------------------------------------------------------------------------
%	REQUIRED PACKAGES AND  CONFIGURATIONS
%------------------------------------------------------------------------------

% CONFIGURATIONS
\usepackage{parskip} % For paragraph layout
\usepackage{setspace} % For using single or double spacing
\usepackage{emptypage} % To insert empty pages
\usepackage{multicol} % To write in multiple columns (executive summary)
\setlength\columnsep{15pt} % Column separation in executive summary
\setlength\parindent{0pt} % Indentation
\raggedbottom

% PACKAGES FOR TITLES
\usepackage{titlesec}
\titlespacing{\section}{0pt}{3.3ex}{2ex}
\titlespacing{\subsection}{0pt}{3.3ex}{1.65ex}
\titlespacing{\subsubsection}{0pt}{3.3ex}{1ex}
\usepackage{color}

% PACKAGES FOR LANGUAGE AND FONT
\usepackage[english]{babel}
\usepackage[utf8]{inputenc}
\usepackage[T1]{fontenc}
\usepackage[11pt]{moresize}

% PACKAGES FOR IMAGES
\usepackage{graphicx}
\usepackage{transparent}
\usepackage{eso-pic}
\usepackage{subfig}
\usepackage{tikz}
\usetikzlibrary{}
\graphicspath{{./Images/}}
\usepackage{caption}
\usepackage{xcolor}
\usepackage{amsthm,thmtools,xcolor}
\usepackage{float}

% STANDARD MATH PACKAGES
\usepackage{amsmath}
\usepackage{amsthm}
\usepackage{amssymb}
\usepackage{amsfonts}
\usepackage{bm}
\usepackage[overload]{empheq}
\usepackage{fix-cm}

% PACKAGES FOR TABLES
\usepackage{tabularx}
\usepackage{longtable}
\usepackage{colortbl}

% PACKAGES FOR ALGORITHMS (PSEUDO-CODE)
\usepackage{algorithm}
\usepackage{algorithmic}

% PACKAGES FOR REFERENCES & BIBLIOGRAPHY
\usepackage[colorlinks=true,linkcolor=black,anchorcolor=black,citecolor=black,filecolor=black,menucolor=black,runcolor=black,urlcolor=black]{hyperref}
\usepackage{cleveref}
\usepackage[square, numbers, sort&compress]{natbib}
\bibliographystyle{abbrvnat}

% OTHER PACKAGES
\usepackage{pdfpages}
\usepackage{afterpage}
\usepackage{fancyhdr}
\fancyhf{}

% Input of configuration file. Do not change config.tex file unless you really know what you are doing.
\input{Configuration_Files/config}

%----------------------------------------------------------------------------
%	BEGIN OF DOCUMENT
%----------------------------------------------------------------------------

\begin{document}

\fancypagestyle{plain}{%
\fancyhf{}
\fancyhead[RO,RE]{\thepage}
\renewcommand{\headrulewidth}{0pt}
\renewcommand{\footrulewidth}{0pt}}

%----------------------------------------------------------------------------
%	TITLE PAGE
%----------------------------------------------------------------------------

\pagestyle{empty}
\frontmatter

\puttitle{
	title=Architectures and AI Moderation Strategies for Multi-Party Speech Conversational Systems,
	name=Salvatore Conte,
	course=Computer Science and Engineering - Ingegneria Informatica, % TODO: verify exact course name
	ID  = 000000,  % TODO: insert matricola
	advisor= Prof. Franca Garzotto,
	coadvisor={}, % TODO: insert co-supervisor if any, otherwise remove this line
	academicyear={2025-26},
}

%----------------------------------------------------------------------------
%	PREAMBLE PAGES: ABSTRACT (EN + IT)
%----------------------------------------------------------------------------
\startpreamble
\setcounter{page}{1}

% ABSTRACT IN ENGLISH
\chapter*{Abstract}
% TODO: write at the very end, once the rest of the thesis is finished.
% Concise summary (single page) of: motivation, research questions,
% architecture and moderation strategies contributed, empirical evaluation
% design and results, relevance and impact.
\\
\\
\textbf{Keywords:} multi-party conversational AI, speech-based systems, AI moderation, turn-taking, empirical evaluation % TODO: refine

% ABSTRACT IN ITALIAN
\chapter*{Abstract in lingua italiana}
% TODO: traduzione dell'abstract inglese, scritta per ultima.
\\
\\
\textbf{Parole chiave:} intelligenza artificiale conversazionale multi-utente, sistemi vocali, moderazione AI, gestione dei turni, valutazione empirica % TODO

%----------------------------------------------------------------------------
%	LIST OF CONTENTS
%----------------------------------------------------------------------------

\thispagestyle{empty}
\tableofcontents
\thispagestyle{empty}
\cleardoublepage

%-------------------------------------------------------------------------
%	THESIS MAIN TEXT
%-------------------------------------------------------------------------

\addtocontents{toc}{\vspace{2em}}
\mainmatter

\chapter{Introduction}
\label{ch:introduction}

% TODO: write last (introductions are easier to write once the rest is in place).
% Structure mirrors the host lab's thesis format:
%   1.1 Context     -> background and motivation
%   1.2 Goals       -> objectives, research questions, contributions
%   1.3 Structure   -> chapter-by-chapter roadmap

This chapter introduces the context, motivations, and objectives of the
thesis. Section~\ref{sec:intro-context} frames the background of
multi-party speech-based conversational systems and AI moderation.
Section~\ref{sec:intro-goals} states the goals of the thesis, the
research questions, and summarises the contributions.
Section~\ref{sec:intro-structure} provides a roadmap of the remaining
chapters.

\section{Context}
\label{sec:intro-context}
% TODO: background on multi-party voice systems, motivation for AI moderation,
% gap that this work addresses (link to Chapter 2).

\section{Goals of the Thesis}
\label{sec:intro-goals}
% TODO: state goals + research questions + summary of contributions.
% RQ1 — How can a speech-based conversational system be architected to support
%       multi-party group discussions with AI moderation?
% RQ2 — What strategies can the AI moderator adopt to intervene effectively
%       without compromising the dynamics of the group?
% RQ3 — How does the presence/absence of the AI moderator affect group
%       performance, participation balance, and user experience in a
%       collaborative survival task?

\section{Structure of the Thesis}
\label{sec:intro-structure}
% TODO: one short paragraph per chapter that follows.

\chapter{State Of The Art}
\label{ch:state-of-the-art}

This chapter surveys the research landscape relevant to the design of speech-based conversational agents for multi-party group interactions. Section~\ref{sec:speech-architectures} reviews the architectural families for speech-based conversational systems (cascaded, half-cascade, end-to-end) and compares their trade-offs. Section~\ref{sec:turntaking} examines turn-taking, from foundational linguistic models to recent machine-learning approaches, with attention to the multi-party setting. Section~\ref{sec:vad-vap} covers Voice Activity Detection and Voice Activity Projection, the low-level perceptual building blocks of turn management. Section~\ref{sec:multiparty-ca} presents multi-party conversational AI as a field, reviewing the main taxonomies and the technical obstacles that have kept the research largely text-based. Section~\ref{sec:ai-roles} surveys the spectrum of roles that AI agents can play in group conversation, from peer participants to full moderators. Section~\ref{sec:group-decision-making} connects this body of work to the broader literature on group decision-making and to the survival-task paradigm used to study it empirically. Section~\ref{sec:discussion} closes the chapter by drawing together the recurring gaps that emerge across these threads.

% ============================================================
%  \section: Speech-Based Conversational Architectures
%  Parent: \section{State of the Art}
% ============================================================

\section{Speech-Based Conversational Architectures}
\label{sec:speech-architectures}

Speech-based conversational agents have to map an incoming acoustic signal
to a spoken response. Ji et al., in the WavChat survey~\cite{ji2024wavechat},
group the architectures used for this task into three families: cascaded,
half-cascade (or semi-cascaded), and end-to-end. The three differ in how
much of the pipeline (speech understanding, language reasoning, speech
synthesis) lives inside a single neural model. The following sections
review each family and discuss the trade-offs.

% ----------------------------------------------------------
\subsection{Cascaded Pipeline (STT--LLM--TTS)}
\label{subsec:cascaded}

The cascaded approach decomposes the problem into three independently
developed stages: an Automatic Speech Recognition (ASR) module transcribes
the user's utterance into text, a Large Language Model (LLM) produces a
textual response, and a Text-to-Speech (TTS) engine renders that response
as audio. Because each component communicates through a well-defined
textual interface, the pipeline is modular: individual stages can be
upgraded, replaced, or scaled without retraining the others.

This architecture underpins the majority of commercially deployed voice
assistants, including the speech pipelines offered by Azure, Google Cloud,
and Amazon Web Services. In the research literature,
X-Talk~\cite{lin2025xtalk} proposes a cascaded dialogue framework targeting
low-latency interaction, while ESPnet-SDS~\cite{li2025espnetsds} provides
an open-source toolkit that implements cascaded spoken-dialogue systems
with pluggable ASR, LLM, and TTS back-ends.

The main advantages of the cascaded pipeline are its maturity and
interpretability. The intermediate text representation makes logging,
debugging, and content moderation straightforward, which is essential in
enterprise and safety-critical settings. Training costs are also low,
since each stage can rely on pre-trained, off-the-shelf models.

The main drawback is latency. The sequential execution of three distinct
inference steps introduces cumulative delays in the range of two to five
seconds, far above the approximately 200\,ms inter-turn gap that
characterises natural human conversation~\cite{heldner2010pauses}.
Streaming reduces this latency in practice. The ASR processes audio while
the user is still speaking, the LLM emits tokens as it generates them, and
the TTS plays each chunk before the full response is ready.
Whisper-Streaming~\cite{machacek2023whisper} reports sub-second incremental
latency on long-form speech using this approach. Errors also compound
across stages: an ASR misrecognition propagates into the LLM prompt and
may trigger an irrelevant response. Finally, the text bottleneck discards
paralinguistic information such as prosody, emotion, and hesitation
markers, which are often crucial for pragmatic interpretation, and prevents
native support for full-duplex interaction.

% ----------------------------------------------------------
\subsection{Half-Cascade (Semi-Cascaded) Architectures}
\label{subsec:half-cascade}

Half-cascade designs fuse two of the three stages and keep the third as
a separate component. The WavChat taxonomy~\cite{ji2024wavechat}
distinguishes three variants, each removing a different inter-stage
boundary.

\emph{Speech-In fusion} (Variant~A) removes the ASR stage: the LLM
ingests audio features directly through a learnable encoder and projection
layer, preserving paralinguistic cues that text transcription would
discard (e.g., SALMONN~\cite{tang2024salmonn},
Freeze-Omni~\cite{wang2024freezeomni}). \emph{Speech-Out fusion}
(Variant~B) absorbs the TTS stage into the LLM, which generates discrete
speech tokens alongside text and so reduces output latency (e.g.,
LLaMA-Omni~\cite{fang2024llamaomni}). \emph{Inner Monologue} (Variant~C)
keeps text as an internal reasoning scaffold while producing speech output
directly, combining coherent reasoning with expressive synthesis (e.g.,
SpeechGPT~\cite{zhang2023speechgpt}).

These architectures cut latency and use richer multimodal representations
than cascaded pipelines, but at a cost. Aligning audio and text
representation spaces needs carefully staged training, and the model's
language abilities can degrade after fine-tuning on speech data, an effect
known as reasoning degradation~\cite{chen2025urobench}. Speech quality
also tends to fall behind dedicated TTS systems, and none of these models
has been evaluated beyond dyadic interaction.

% ----------------------------------------------------------
\subsection{End-to-End Speech-to-Speech Models}
\label{subsec:e2e}

End-to-end architectures collapse the entire spoken-dialogue pipeline into
a single neural model that maps audio input directly to audio output. A
key enabler is the family of \emph{neural audio codecs} (SoundStream,
EnCodec, Mimi), which convert continuous waveforms into compact sequences
of discrete tokens suitable for autoregressive modelling.

Moshi~\cite{defossez2024moshi} was the first publicly demonstrated
full-duplex speech-to-speech model, with a mouth-to-ear latency of about
200\,ms. Commercial systems such as GPT-4o (OpenAI), Gemini~2.5 (Google),
Qwen3-Omni (Alibaba), and Kimi Audio (Moonshot AI) have since adopted
end-to-end or near-end-to-end designs, although their architectural
details remain mostly proprietary. In the academic domain, dGSLM extends
the Generative Spoken Language Model to dialogue by jointly modelling two
audio channels.

End-to-end models have a clear appeal: removing the intermediate
representations cuts latency and preserves paralinguistic cues (pitch
contour, speaking rate, emotional colouring) that cascaded systems
discard. They also support full-duplex interaction natively, since the
model can listen and speak at the same time.

Current end-to-end systems still have important limitations.
URO-Bench~\cite{chen2025urobench} shows that, on reasoning-intensive
tasks, end-to-end speech models perform substantially worse than their
text-only LLM counterparts, suggesting that the speech modality introduces
a reasoning bottleneck. Training costs are high, since they need large
volumes of paired speech-to-speech data. Reliability is also a concern:
Lin et al.~\cite{lin2025xtalk} observe that end-to-end speech models are
``often unsuitable for deployment'' in production settings due to
hallucination, inconsistent output quality, and limited controllability.

% ----------------------------------------------------------
% Closing paragraph: ties the architecture analysis to the multi-party
% setting that motivates AIutami's cascaded choice. Not a separate
% subsection on purpose, to avoid the "research gap" framing — the
% chapter-level gap discussion lives in the final Discussion section.

Every multi-party speech-based AI system reported in the literature is
built on a cascaded pipeline. The ARI robot~\cite{addlesee2024ari} and
the Furhat platform~\cite{axelsson2025furhat} both use cascaded STT, LLM,
and TTS to handle speaker identification, overlapping speech, and group
turn-taking. End-to-end and half-cascade models have so far only been
tested in dyadic settings. This is an empirical observation, not a
theoretical limit: nothing prevents a neural speech-to-speech model from
supporting groups in principle, but no published work shows it actually
working. A multi-party system built today therefore commits to the
cascaded pattern, and the interesting design choices sit at the layers
above the pipeline: turn-taking, addressee detection, and moderation.
These are the focus of the rest of this chapter.


% ==============================================================
%  Turn-Taking in Spoken Dialogue  &  VAD / VAP
%  To be \input inside \section{State of the Art}
% ==============================================================

\section{Turn-Taking in Spoken Dialogue}
\label{sec:turntaking}

Turn-taking---the mechanism by which interlocutors coordinate who speaks when---is one of the most fundamental structures of human conversation. Any spoken dialogue system that aspires to natural interaction must address this coordination problem, yet doing so computationally remains a significant challenge, particularly in multi-party settings.

\subsection{Foundational Framework}

The seminal work of Sacks, Schegloff, and Jefferson~\cite{sacks1974turntaking} established the analytical framework that still underpins most computational approaches to turn-taking. Their model identifies \emph{Transition Relevance Places} (TRPs)---points in an utterance at which a change of speaker becomes appropriate---and proposes three ordered rules governing speaker selection: (1)~the current speaker selects the next speaker, (2)~a non-current speaker self-selects, or (3)~the current speaker may continue. Empirical work by Heldner and Edlund~\cite{heldner2010pauses} later quantified the remarkable efficiency of this system, showing that inter-turn gaps in human conversation average approximately 200\,ms---a latency so short that it implies listeners begin planning their response well before the current speaker finishes.

This observation has profound implications for dialogue system design: a spoken agent that waits for silence and then begins processing the input will inevitably feel unresponsive compared to a human interlocutor.

\subsection{Computational Approaches}

Four broad families of computational approaches to turn-taking can be distinguished, representing an evolution from simple heuristics to sophisticated predictive models.

\paragraph{Rule-Based and Threshold-Based Methods.}
The most common approach in commercial dialogue systems relies on fixed silence thresholds: once no speech energy is detected for a predetermined interval---typically between 500 and 700\,ms---the system assumes the user has finished speaking. This technique, known as \emph{endpointing}, is employed by most major voice assistants including Amazon Alexa and Apple Siri. While computationally simple and robust, threshold-based endpointing suffers from an inherent trade-off: short thresholds lead to premature cutoffs during within-turn pauses (e.g., hesitations, breathing, or mid-sentence planning), whereas long thresholds introduce perceptible latency that disrupts conversational flow.

\paragraph{Prosody-Based Methods.}
A more linguistically informed family of approaches exploits the prosodic cues that human speakers use to signal turn boundaries. Gravano and Hirschberg~\cite{gravano2011turncues} conducted a comprehensive analysis of turn-yielding cues in the Columbia Games Corpus, identifying six acoustic features---including falling pitch, final syllable lengthening, reduced intensity, and specific discourse markers---that co-occur with turn boundaries. The more such cues are present simultaneously, the higher the probability that the interlocutor will take the floor. In a related line of work, Ward and Tsukahara~\cite{ward2000backchannel} demonstrated that low pitch regions in the speaker's signal reliably predict opportunities for listener backchannels, providing a basis for systems that can produce appropriately timed feedback signals.

\paragraph{Machine Learning Methods.}
Raux and Eskenazi~\cite{raux2009endpointing} introduced the first machine-learning-based endpointing system, training a decision tree on a combination of prosodic, lexical, and dialogue-level features to classify pause points as either turn-internal or turn-final. Their system, deployed in the CMU Let's Go bus information system, significantly reduced both false endpoint detections and average response latency compared to a fixed-threshold baseline. Skantze~\cite{skantze2017turntaking} subsequently reframed turn-taking as a \emph{continuous prediction problem}: rather than making binary decisions at discrete moments, a system should continuously estimate the probability that the current speaker will yield the floor, that a backchannel is appropriate, or that the system should hold its turn. Using recurrent neural networks (LSTMs), this approach enabled more fluid and anticipatory behaviour. Meena et al.~\cite{meena2014turntaking} further explored data-driven approaches in a spoken dialogue context, demonstrating the effectiveness of combining prosodic and contextual features for turn-taking decisions. Ekstedt and Skantze~\cite{ekstedt2020turngpt} later proposed \emph{TurnGPT}, a Transformer-based model that predicts turn shifts from linguistic content alone, showing that a language model fine-tuned on dialogue data can learn implicit turn-yielding patterns even without access to acoustic features.

\paragraph{LLM-Based Methods.}
The most recent line of work leverages large language models directly for turn-taking decisions. Addlesee~\cite{addlesee2024ari} proposed using an LLM prompt to assess whether an incoming user utterance constitutes a complete conversational contribution, thereby deciding whether to respond or to wait for additional input. Houde et al.~\cite{houde2025koala} developed \emph{Koala}, a system in which the LLM assigns a self-score to its confidence that the user has finished speaking, enabling a soft and adjustable endpointing mechanism. These approaches benefit from the world knowledge and pragmatic competence encoded in large pre-trained models, though they introduce additional latency due to inference time and depend on the LLM's capacity to reason reliably about conversational structure.

\subsection{Key Concepts}

Table~\ref{tab:turntaking-concepts} summarises the core terminology used across the turn-taking literature.

\begin{table}[ht]
  \centering
  \caption{Key turn-taking concepts and definitions.}
  \label{tab:turntaking-concepts}
  \begin{tabularx}{\textwidth}{@{} l X @{}}
    \toprule
    \textbf{Concept} & \textbf{Definition} \\
    \midrule
    Turn-yielding   & Behaviour indicating the current speaker intends to relinquish the floor (e.g., falling pitch, syntactic completion, gaze shift). \\
    Turn-holding     & Signals that the current speaker wishes to retain the floor despite a pause (e.g., filled pauses, rising intonation, averted gaze). \\
    Backchannel      & Short listener feedback (e.g., ``mhm'', ``yeah'') that signals attention without claiming the floor. \\
    Barge-in         & A listener begins speaking while the current speaker is still talking, explicitly claiming the turn. \\
    Floor            & The right or opportunity to speak at a given moment; in multi-party settings, multiple notions of floor may coexist. \\
    Overlap          & Simultaneous speech by two or more participants, which may be collaborative (e.g., backchannels) or competitive (e.g., interruptions). \\
    TRP              & \emph{Transition Relevance Place}: a point at which speaker change becomes appropriate, typically at syntactic or pragmatic boundaries. \\
    Endpointing      & The computational task of detecting that a user has finished their turn, enabling the system to begin response generation. \\
    \bottomrule
  \end{tabularx}
\end{table}

\subsection{Multi-Party Turn-Taking}

While the above approaches have been developed and evaluated primarily in dyadic (two-party) settings, multi-party conversations introduce qualitatively different challenges. The ``who's next?'' problem becomes substantially harder when more than two participants are present: the system must not only determine \emph{whether} a turn transition is occurring but also \emph{who} should speak next. In face-to-face interaction, gaze direction serves as the primary mechanism for addressee selection and next-speaker nomination, but this channel is often unavailable or unreliable in voice-only or remote settings. Additional complications include overlapping speech from multiple participants, the need for speaker diarisation to attribute utterances to individuals, and the emergence of complex social dynamics such as schisms (parallel sub-conversations within a group), dominance patterns, and unequal participation.

Several turn allocation mechanisms have been explored in the literature. \emph{Explicit nomination} requires the current speaker to name the next participant verbally. \emph{Gaze-based allocation} exploits visual attention as an implicit selection signal; Johansson et al.~\cite{johansson2013gaze} demonstrated that a robot capable of interpreting and producing gaze cues could manage multi-party turn transitions more naturally. \emph{Self-selection races}, as described in the original model of Sacks et al.~\cite{sacks1974turntaking}, allow any participant to claim the floor at a TRP, with the first speaker to begin typically gaining precedence. Finally, \emph{moderator-mediated} allocation delegates turn management to a designated party---whether human or artificial---that grants or denies speaking rights according to some policy.

From a computational modelling perspective, Bohus and Horvitz~\cite{bohus2009engagement,bohus2011multiparty} developed influential engagement models for multi-party interaction with embodied agents, addressing the problems of engagement recognition, maintenance, and directed-action selection in settings with multiple co-present users.

% ==============================================================
\section{Voice Activity Detection and Projection}
\label{sec:vad-vap}

Voice Activity Detection (VAD) and Voice Activity Projection (VAP) represent two ends of a temporal spectrum in the analysis of spoken interaction: the former determines whether speech is occurring \emph{now}, while the latter predicts whether speech will occur in the \emph{near future}. Together, they provide the low-level perceptual foundation upon which higher-level dialogue phenomena---such as turn-taking, backchannel production, and barge-in detection---can be built.

\subsection{Voice Activity Detection}

Voice Activity Detection (VAD) is the task of classifying, at each time step, whether speech is present in an audio signal.  Early systems relied on hand-crafted acoustic features and lightweight statistical models; the widely adopted WebRTC VAD exemplifies this generation, offering low-latency frame-level decisions at minimal computational cost, albeit with limited noise robustness.  Modern neural VAD systems---such as \emph{Silero VAD}, the \emph{pyannote-audio} pipeline, and Google's \emph{Personal VAD}---achieve substantially higher accuracy in adverse conditions while remaining efficient enough for real-time deployment.

In spoken dialogue, VAD fulfils three primary roles: it \emph{gates} the ASR module so that only speech-containing segments are transcribed, it provides the foundational signal for \emph{endpointing} (detecting when a speaker has finished), and it enables \emph{barge-in detection} during system output.  Crucially, however, VAD determines only \emph{that} someone is speaking, not \emph{who} is speaking (which requires speaker diarisation) nor whether the detected speech constitutes a complete turn or merely a backchannel (which requires higher-level turn-taking prediction).

\subsection{Voice Activity Projection}

While VAD operates reactively on the present signal, Voice Activity Projection~(VAP), introduced by Ekstedt and Skantze~\cite{ekstedt2022vap}, takes a fundamentally predictive stance: given the audio observed so far, the model forecasts the \emph{future} voice activity of each participant over a horizon of 0.5 to 2 seconds. This predictive capability enables a dialogue system to anticipate turn transitions, backchannel opportunities, and overlap onsets before they occur, rather than merely detecting them after the fact.

\paragraph{Architecture and Training.}
The VAP model receives stereo audio (one channel per speaker in a dyadic conversation) and extracts frame-level features using a self-supervised speech representation model such as CPC or wav2vec~2.0. These features are processed by a Transformer encoder, which outputs a probability distribution over future voice activity states at each time step. The future is discretised into four temporal bins per speaker; with two speakers, this yields $2^{4 \times 2} = 256$ possible activity state combinations. The model is trained in a \emph{self-supervised} fashion: the actual future voice activity observed in the data serves directly as the training signal, eliminating the need for any manual annotation of turn-taking events.

\paragraph{Emergent Capabilities.}
A remarkable property of the VAP framework is that complex turn-taking behaviours emerge implicitly from the simple objective of predicting future voice activity. Without ever being explicitly trained on turn-taking labels, the model learns to recognise turn-yielding cues, identify appropriate moments for backchannels, and distinguish turn-holds from turn-final silences. This suggests that much of the information required for competent turn-taking behaviour is already encoded in the acoustic signal and can be recovered through prediction of a straightforward observable quantity. TurnGPT~\cite{ekstedt2020turngpt} had previously demonstrated an analogous phenomenon in the text modality: a language model trained on dialogue transcripts implicitly learned turn-shift patterns from linguistic content alone.

\subsection{The Multi-Party Scalability Problem}

The current VAP architecture is designed for dyadic interaction and assumes a stereo audio input with exactly two speaker channels. Extending it to multi-party settings poses a fundamental scalability challenge. With $N$ participants and four temporal bins per speaker, the number of possible future voice activity states grows exponentially as $2^{4N}$. For a conversation with just five participants, this already yields over one million states ($2^{20} = 1{,}048{,}576$), making a direct extension of the output representation infeasible.

An alternative approach---processing all $\binom{N}{2} = \frac{N(N-1)}{2}$ pairwise speaker combinations independently---does not scale either, as it grows quadratically with the number of participants and fails to capture group-level dynamics that are not reducible to pairwise interactions. Malik, Atieh, and Skantze~\cite{malik2024vap} have explicitly identified the extension of VAP to multi-party conversation as an open research problem, exploring initial strategies such as dimensionality reduction of the output space, hierarchical prediction schemes, and attention-based architectures that can accommodate a variable number of input channels. This remains an active area of investigation with no established solution.



% ============================================================
%  \section: Multi-Party Conversational AI
%  \section: Research Gaps
%  \section: Positioning of AIutami
%  Parent: \section{State of the Art}
% ============================================================

\section{Multi-Party Conversational AI}
\label{sec:multiparty-ca}

The conversational agents in widespread use today, including Siri, Alexa, ChatGPT, and GPT-4o voice mode, are designed for dyadic interaction: one human speaks to one artificial interlocutor, alternating turns in a strict A-B-A pattern. Real human communication, on the other hand, is mostly group-based. Meetings, classroom discussions, therapy sessions, and brainstorming workshops all involve several participants whose contributions interleave, overlap, and address different subsets of the group.

The move from dyadic to multi-party conversation is a qualitative paradigm shift, not an incremental extension~\cite{gu2022mpc,zheng2022polyadic,traum2004issues}. In a dyadic exchange the information flow is sequential and predictable. In a multi-party setting it becomes a graph: any utterance can come from any participant and can be addressed to any individual, any subset, or the whole group, which gives rise to parallel sub-conversations, shifting coalitions, and complex turn-allocation patterns~\cite{gu2022mpc}. These structural differences propagate to every layer of a conversational AI system, from speaker identification and turn management to response generation and addressee resolution, and they call for dedicated modelling strategies that have no counterpart in dyadic research.

The rest of this section reviews three taxonomies that capture the computational, design, and user-experience dimensions of multi-party conversational AI, and then discusses the technical obstacles that explain why most multi-party research has so far stayed text-based.

\paragraph{Gu et al. (2022): WHO Says WHAT to WHOM.}
Gu et al.~\cite{gu2022mpc} present an IJCAI survey that decomposes multi-party conversation into three interrelated computational problems. \emph{Speaker Modelling} (WHO) is about identifying and predicting the current and next speaker; solutions range from Conditional Random Fields and LSTMs to Transformer-based turn-taking predictors. \emph{Utterance Modelling} (WHAT) addresses what the agent should say, with retrieval-based methods such as TopicBERT alongside generation-based approaches like HeterMPC and, more recently, LLM prompting. \emph{Addressee Modelling} (WHOM) tackles the problem of who an utterance is directed at: addressee recognition, dialogue disentanglement, and the disambiguation of shared-channel messages. Gu et al.\ note that the three problems are usually studied in isolation, but they are deeply interdependent in practice: knowing who speaks conditions the interpretation of what is said, which in turn constrains the inference of whom it addresses. A key limitation of their survey is that it covers only text-based multi-party conversation. Extending it to speech adds further complexity, including ASR errors, overlapping speech, prosodic cues, and real-time latency constraints, none of which the taxonomy addresses.

\paragraph{Houde et al. (2025): designing the agent's participation.}
Houde et al.~\cite{houde2025koala} propose an HCI-oriented framework, developed at IUI~2025, for specifying and controlling the behaviour of an AI agent in group conversations. The framework is organised along two orthogonal axes. The first axis covers \emph{what to control}: WHEN governs the triggers, filters, and frequency of agent contributions (only when addressed, when a self-assessed relevance score exceeds a threshold, or at a fixed cadence); WHAT governs the content, style, and modality of contributions (conservative vs.\ creative ideas, formal vs.\ casual tone); WHERE governs the channel used by the agent (a shared group channel, a threaded reply, or a private message). The second axis covers \emph{how to control it}: SPECIFY describes the mechanism for configuring the agent's behaviour (a UI panel, natural-language instructions, role-based presets); ACCESS describes who is authorised to adjust the configuration (an administrator, any participant, or a democratic consensus); IMPLEMENT describes how the configuration is enforced at the system level (system prompt, external logic, or a hybrid). This design taxonomy is the complement of Gu et al.'s decomposition: where the latter describes the technical sub-problems of multi-party conversation, Houde et al.'s framework describes the design space for governing the agent's participation.

\paragraph{Zheng et al. (2022): challenges and design properties.}
Zheng et al.~\cite{zheng2022polyadic} conduct a CHI literature review of 36 papers on polyadic (multi-party) human-AI interaction, distilling four recurrent challenges. \emph{Inefficient communication} covers topic drift, lack of structure, and difficulty reaching consensus. \emph{Lack of engagement} covers unequal participation and passive members. \emph{Relationship-maintenance barriers} include insufficient emotional awareness and difficulty regulating group affect. The fourth, the \emph{need for building connections}, covers ice-breaking, common-ground establishment, and cross-cultural issues. From these challenges they derive three design properties for a polyadic conversational agent. The agent should be \emph{Visible}: participants must be aware that it is there and what its role is. It should be \emph{Ignorable}: its interventions must be non-intrusive, since too many or too long become counterproductive. And it should be \emph{Accountable}: it must be clear who it is responding to and who controls its behaviour. These three properties impose concrete constraints on system design. In text-based environments, visibility can be achieved through distinct usernames or badges, ignorability through message threading, and accountability through explicit @-mentions. In speech-based environments the same properties take on different meanings: the agent's voice is visible by default since it occupies a shared audio channel, it is hard to ignore since one cannot scroll past a spoken utterance, and it is hard to make accountable without explicit turn-assignment mechanisms.

The amount of speech-based multi-party research is small, and the reason is technical. Addlesee et al.~\cite{addlesee2020asr} provide a systematic evaluation of commercial ASR and speaker-diarisation services in dyadic and multi-party settings, and report performance drops that are often too severe for practical use. Speaker diarisation, the task of attributing each audio segment to the correct speaker, has Diarisation Error Rates between 49\% and 67\% in multi-party settings on commercial platforms, and Google's service fails entirely when more than three speakers are present. Overlapping speech, which is natural whenever speakers compete for the floor or produce backchannels, raises Word Error Rates by about 20\%. Incremental ASR, needed for responsive turn-taking, gives a sixfold increase in WER over batch transcription (33\% versus 5\%). Disfluencies such as filled pauses, self-corrections, and edit terms add further variability across services. These compound difficulties explain why the multi-party conversational AI research surveyed by Gu et al.~\cite{gu2022mpc}, Zheng et al.~\cite{zheng2022polyadic}, and Houde et al.~\cite{houde2025koala} has stayed almost exclusively in the text modality.

% ============================================================
\section{AI Agents and their Roles in Group Conversation}
\label{sec:ai-roles}

The literature on group conversation with AI spans a spectrum of agent roles that mirrors a long-standing distinction in group decision research. At one end of the spectrum is the \emph{facilitator}, defined by Schwarz~\cite{schwarz2017facilitator} as an agent that is substantively neutral, has no authority over the decision itself, and intervenes to help the group improve how it identifies and solves problems. At the other end is the \emph{moderator}, a more directive role that enforces structure, allocates turns, and steers the discussion toward a goal. The systems reviewed below populate this spectrum and show what each role implies in terms of design, turn management, and evaluation.

At the most peer-like end of the spectrum the AI sits alongside human members and contributes ideas of its own. Koala~\cite{houde2025koala} is a Slack-based chatbot deployed in three-person human brainstorming groups at IBM. Houde et al.\ evaluate two variants: in the \emph{reactive} variant Koala contributes only when explicitly addressed via @-mention, while in the \emph{proactive} variant it uses a self-scoring LLM mechanism that estimates the value of a possible contribution and posts when the score exceeds a configurable threshold. The proactive variant generated 73\% of all ideas in the AI-augmented condition but only 33\% of the top-rated ones, and 72\% of participants preferred the reactive variant, describing the proactive agent as ``a pedantic student who wouldn't create space for others''. A follow-up study, Koala~II, added a user-tunable proactivity threshold: no group reverted to fully reactive mode, and the perceived usefulness of the controls was rated 4.46 out of 5. The central finding is that agent proactivity should be neither binary nor static, and that users want to adjust it dynamically during a session.

Closer to the centre of the spectrum, the AI takes a supporting role: it responds when addressed but does not steer the discussion. Two recent robot-based systems illustrate this pattern. The ARI robot~\cite{addlesee2024ari} is deployed in a hospital memory clinic, where it interacts with two human participants (a patient and a companion) via speech. ARI uses a cascaded pipeline with gaze-based addressee detection through an LLM classifier, reaching 85.4\% accuracy compared to 53.4\% with text alone. The system also implements incremental Clarification Requests (iCR) for incomplete utterances, which are common with patients experiencing cognitive decline, and a grounding technique that attributes information to a fictitious persona, reducing hallucination by about 10 percentage points across annotated responses. The Furhat robot~\cite{axelsson2025furhat} is a social robot that engages in open-ended conversation with two human participants. Its pipeline processes voice direction-of-arrival, Azure STT, speaker diarisation, and face recognition, then forwards the result to GPT-3.5 for response generation. Addressee accuracy reaches 92.6\% in parallel settings, drops to 79.3\% in true group conversation, and falls to just 26.8\% with audio alone on concurrent speakers, confirming that audio is not enough to resolve addressee ambiguity. Both robots stay close to dyadic interaction with two human participants, and neither takes a moderating role.

Further along the spectrum, the AI acts as a facilitator: it supports the social process of the group without dominating it, often in emotionally sensitive contexts. The Empathic Co-facilitator~\cite{adikari2022empathic} is a text-based AI co-facilitator for online cancer-support groups (Cancer Chat Canada), analysed across approximately 120\,000 conversational turns. It runs alongside a human therapist, monitors group dynamics, detects emotions using Plutchik's eight-emotion taxonomy via Word2Vec and GloVe embeddings, and models emotion-state transitions through Markov chains. When the group's negativity score exceeds a threshold, the system either generates an empathic response from pre-defined templates or alerts the therapist. The system pre-dates the LLM era, so its language capabilities are limited to template-based responses, and it operates exclusively in asynchronous text. A more recent and directly experimental study by Gao et al.~\cite{gao2025facilitation} compares AI versus human facilitation in peer-support groups for exam anxiety, with 112 participants split between AI-facilitated and human-facilitated discussions. Across all measures, including emotional connection, perceived facilitation effectiveness, perceived utility, and intention to use the system again, human facilitators outperformed the AI. This is one of the few direct experimental comparisons of the two facilitation conditions, and it suggests that AI moderation is not yet a drop-in replacement for human-led group support.

At the most directive end of the spectrum, the AI acts as a full moderator: it allocates turns, enforces structure, and steers the discussion toward a goal. The historical precursor is VIRTUAL HUMAN~\cite{wahlster2023virtual}, a multimodal sports-quiz system from 2006 featuring two human contestants and three virtual agents (one moderator and two domain experts). Dialogue management relied on symbolic rules rather than statistical or neural models. While the system shows that the idea of an AI moderator in multi-party conversation has a two-decade lineage, its pre-LLM architecture makes it very different from contemporary approaches in both flexibility and scalability. A more recent and methodologically closer precedent is the LLM-powered moderator studied by Dubini~\cite{dubini2024multiuser} in a text-based three-person Murder Mystery task adapted from Stasser and Stewart's hidden-profile paradigm. Across 20 sessions recruited on Prolific (10 without moderator, 10 with the LLM moderator), the work reports that the moderator subtly influenced participation dynamics and the focus of the discussion but had limited measurable effect on decision accuracy on the small sample. Dubini framed the system as a baseline rather than a finished design, leaving open the question of how to make AI moderation effective in this kind of task.

The systems above evaluate AI in group conversation on very different tasks and metrics. Koala measures brainstorming productivity and user preference, ARI and Furhat measure addressee accuracy, Empathic and Gao et al.\ measure emotional connection and perceived effectiveness, and Virtual Human is a proof of concept without a controlled evaluation. The choice of task and metric shapes what we can learn about an AI agent in a group. A particularly useful family of tasks for measuring objective group performance, with a long history in social and organisational psychology, are the survival decision-making tasks reviewed in the next section.



% ==============================================================
%  Group Decision-Making and Survival Tasks
% ==============================================================

\section{Group Decision-Making and Survival Tasks}
\label{sec:group-decision-making}

The systems reviewed in Section~\ref{sec:ai-roles} place an AI agent inside group conversations, but the question of what it means for a group to decide well belongs to a different field: group decision-making, studied for decades in social and organisational psychology. Understanding this literature is needed both to frame the problem and to recognise the experimental paradigms the discipline has developed to study it.

The starting point is Steiner's \emph{Group Process and Productivity}~\cite{steiner1972group}, which introduces the notion of \emph{process loss}: a group typically produces less than the sum of its parts, because some individual effort is dissipated in coordination, motivational problems, and inefficient communication. Hill's meta-analysis of decades of comparative studies~\cite{hill1982group} confirms that groups do not systematically beat the best individual: they do under specific conditions and on specific kinds of tasks, but not in general. Stasser and Stewart~\cite{stasser1992hidden} add a specific mechanism known as the \emph{hidden profile} paradigm: when critical information is unevenly distributed among members, so that each one has a unique piece and only by pooling the pieces can the group reach the correct solution, discussion tends to revolve around the information already shared rather than surfacing what is unique, and groups end up making poor decisions. This is the structure of the Murder Mystery task used by Dubini in Section~\ref{sec:ai-roles}. The overall picture is that group performance is contingent: it depends on the task, on the structure of the information, and on the discussion process, which leaves room for procedural guidance, human or artificial, to improve it.

To study these phenomena empirically the discipline has developed \emph{survival decision tasks}: scenarios in which a group has to rank a list of items by survival priority, first individually and then as a group, with performance measured as the deviation from a reference ranking produced by an expert. The original task is NASA Moon Survival, developed by Hall in the early 1960s and formalised in the experimental study by Hall and Watson~\cite{hall1970normative}: fifteen items to rank for survival after an emergency landing on the Moon, with an expert ranking supplied by NASA. Hall and Watson report a result that is central to the present discussion: groups receiving explicit consensus-seeking instructions (look for disagreement, avoid majority voting, do not yield just to avoid conflict) produce significantly better decisions than groups without instructions, and beat their most competent member more often (strong synergy). Nemiroff and Pasmore~\cite{nemiroff1975lostatsea} propose a sibling task, Lost at Sea, with the same structure but a maritime scenario. The paradigm is still active in contemporary research: Hamada and colleagues~\cite{hamada2020wisdom} use NASA Moon Survival to compare group deliberation with wisdom-of-crowds aggregation across twenty-five face-to-face groups, and Hémon and colleagues~\cite{hemon2024constructive} use a survival task in videoconference settings to show that structured instructions based on constructive controversy boost synergy in online groups as well.


% ============================================================
%  Discussion (final section of the State of the Art)
% ============================================================

\section{Discussion}
\label{sec:discussion}

The previous sections have reviewed the architectures available for building speech-based conversational systems, the turn-taking mechanisms that coordinate who speaks when, the VAD and VAP layers that underpin them, the small but growing literature on multi-party conversational AI, the spectrum of roles that an AI agent can play inside a group, and the work on group decision-making with its survival-task paradigm. Taken together, this body of work points to a few recurring gaps that are visible across more than one of these threads.

The first is a modality gap. Most of the work on AI agents in group conversation is text-based. The surveys by Gu et al.~\cite{gu2022mpc}, Zheng et al.~\cite{zheng2022polyadic}, and Houde et al.~\cite{houde2025koala} consider almost exclusively chat-style interactions, and the most recent study of AI moderation in a decision-making task, by Dubini~\cite{dubini2024multiuser}, also runs in chat. The systems that operate in speech are restricted to dyadic settings with two human participants and a gaze channel, such as ARI~\cite{addlesee2024ari} and Furhat~\cite{axelsson2025furhat}. The reason is technical: Addlesee et al.~\cite{addlesee2020asr} report that commercial diarisation services have error rates between 49\% and 67\% in multi-party settings, and that overlapping speech adds another 20\% to the word error rate. As a result, audio-only group conversation with more than two human participants is still largely unexplored.

The second gap concerns turn-taking in groups. In dyadic settings, predictive turn-taking based on Voice Activity Projection~\cite{ekstedt2022vap} has produced models that anticipate speaker transitions from acoustic features alone. Extending this to multi-party is an open problem: Malik et al.~\cite{malik2024vap} note that the output space of VAP grows exponentially with the number of speakers, and no current model handles more than two channels reliably. Audio-only addressee detection collapses to about 27\% accuracy with concurrent speakers~\cite{axelsson2025furhat}, and the end-to-end speech-to-speech models such as Moshi~\cite{defossez2024moshi}, LLaMA-Omni~\cite{fang2024llamaomni}, and Freeze-Omni~\cite{wang2024freezeomni} are all designed for one human at a time. A multi-party speech system built today therefore has to commit to a cascaded pipeline with an explicit floor-management protocol, and the design effort shifts up to the layers above the pipeline.

The third gap is on the role of the AI in the group. Koala~\cite{houde2025koala} reports that a proactive AI in a text group was described by participants as ``a pedantic student who wouldn't create space for others''. Gao et al.~\cite{gao2025facilitation} compare an AI facilitator to a human facilitator in peer-support groups and find that the human still wins on every emotional measure. Dubini~\cite{dubini2024multiuser} reports that an LLM moderator in a text-based Murder Mystery task has only a subtle effect on participation dynamics and limited effect on decision accuracy. The group decision-making literature is more encouraging: Hall and Watson~\cite{hall1970normative} show that explicit consensus-seeking instructions improve group performance on survival tasks, and Hémon et al.~\cite{hemon2024constructive} show that the same effect carries over to online groups working through videoconferencing. The combined picture is that structured guidance helps groups, but the right amount and form of structure for an AI moderator in audio-only multi-party settings is not yet established.

A fourth gap is on evaluation. The text-based work on multi-party AI typically measures user preference and subjective satisfaction~\cite{houde2025koala,zheng2022polyadic}. The robotic speech work measures addressee accuracy~\cite{addlesee2024ari,axelsson2025furhat}. Empathic Co-facilitator~\cite{adikari2022empathic} and Gao et al.~\cite{gao2025facilitation} measure emotional connection. None of these works measures group performance in the sense that the group decision-making literature does, with deviation from an expert ranking and synergy with respect to the best individual member. The survival-task paradigm provides such measures, but it has not so far been combined with AI-moderated speech-based interaction in a controlled empirical study.

In summary, AI moderation of group conversation is an active line of research, but it has so far stayed in text, or in dyadic speech with gaze, or in single-task chatbot prototypes. Speech-based multi-party AI moderation, with audio-only operation, more than two human participants, and a quantitative evaluation grounded in the group decision-making literature, is an open space. The chapters that follow describe one attempt to occupy it.


\chapter{Requirements \& Design}
\label{ch:requirements-design}

% Structure mirrors the host lab's thesis format (7 sections):
%   3.1 State Of The Art Informed Requirements
%   3.2 User Centered Requirements
%   3.3 Design Requirements Consolidation
%   3.4 User Experience (UX) Design
%   3.5 User Interface (UI) Design
%   3.6 Moderator Agent Design   (= "Assistant Design" in the host-lab format)
%   3.7 Use Case

This chapter presents the requirements that drove the design of AIutami
and the resulting design choices. The requirements are first derived
from the literature gaps surfaced in Chapter~\ref{ch:state-of-the-art}
(Section~\ref{sec:rd-sota-req}) and from user-centred enquiry with
target participants and the supervising lab
(Section~\ref{sec:rd-user-req}). Section~\ref{sec:rd-consolidation}
consolidates the two streams into a unified requirement list.
Sections~\ref{sec:rd-ux},~\ref{sec:rd-ui},
and~\ref{sec:rd-moderator-agent} describe respectively the user
experience, the user interface, and the moderator agent design.
Section~\ref{sec:rd-use-case} closes the chapter with a representative
use case that walks through a full group session.

\section{State Of The Art Informed Requirements}
\label{sec:rd-sota-req}
% TODO: requirements derived from the gaps surfaced in Chapter 2.
% Examples:
%  - Support N participants in real-time speech (multi-party, not dyadic)
%  - Explicit turn-taking (mitigate multi-party scalability problem)
%  - Pluggable task layer (NASA Moon, Lost at Sea, Murder Mystery, Generic)
%  - Comparable control arm (moderator OFF) for empirical evaluation

\section{User Centered Requirements}
\label{sec:rd-user-req}
% TODO: requirements gathered from informal feedback with the host lab and
% from the pilot participants (tutors). Examples:
%  - Mobile-friendly UI (participants often join from phones)
%  - Robust rejoin after disconnect (mobile networks are unstable)
%  - Visible turn indicator (no doubt about who is speaking)
%  - GDPR-compliant data collection (Italian university study)

\section{Design Requirements Consolidation}
\label{sec:rd-consolidation}
% TODO: merge SOTA-derived + user-derived into a single numbered list
% (R1, R2, ...). Cite each one back to its source section.

\section{User Experience (UX) Design}
\label{sec:rd-ux}
% TODO: describe the user journey:
% Lobby -> Individual ranking (NASA/LAS only) -> Group discussion (ACTIVE)
% -> Group ranking submission -> Conclusion -> Report.
% Mention design choices: avatar of moderator, ready-to-conclude button,
% participation indicators, etc.

\section{User Interface (UI) Design}
\label{sec:rd-ui}
% TODO: describe the screens (lobby, ranking, session, conclusion, report).
% Include wireframes / screenshots. Discuss accessibility choices.

\section{Moderator Agent Design}
\label{sec:rd-moderator-agent}
% TODO: design rationale for the AI moderator:
%  - Goals: neutrality, minimal intrusion, fairness
%  - Trigger taxonomy: NO_PUSH, INACTIVE_VOICE, TIMER_25/30, soft (LLM)
%  - Intervention categories: questioning, summarising, redirecting, time-keeping
%  - Forced conclusion recap
%  - Prompt engineering decisions

\section{Use Case}
\label{sec:rd-use-case}
% TODO: walk through one representative session (e.g. Lost at Sea, 4 participants,
% moderator ON). Step-by-step narrative covering lobby join, intro, individual
% ranking, group discussion with one or two AI interventions, ready-to-conclude
% flow, and report generation.

\chapter{Implementation}
\label{ch:implementation}

% Structure mirrors the host lab's thesis format (3 sections):
%   4.1 Technologies
%   4.2 Software Architecture
%   4.3 UML Flow diagrams

This chapter documents how the design decisions of
Chapter~\ref{ch:requirements-design} were realised in working software.
Section~\ref{sec:impl-technologies} surveys the technology stack.
Section~\ref{sec:impl-architecture} describes the software architecture
of AIutami, with emphasis on the backend split into purpose-specific
Django apps. Section~\ref{sec:impl-uml} presents the UML diagrams that
capture the runtime behaviour of the most relevant subsystems
(turn-taking, audio forwarding, moderation pipeline, session
lifecycle).

\section{Technologies}
\label{sec:impl-technologies}
% TODO. Suggested coverage:
%   - Backend framework: Django + Daphne (ASGI), Channels for WebSockets
%   - Real-time audio: aiortc (WebRTC server-side), STUN/TURN considerations
%   - AI services: OpenAI Realtime (STT), OpenAI gpt-4o-mini (LLM),
%                  OpenAI gpt-4o-mini-tts (TTS)
%   - Persistence: PostgreSQL 16 (sessions, participants, votes, events),
%                  Redis 7 (turn state, ASR cache, moderation state)
%   - Frontend (separate repo): React + Vite + Tailwind, mobile-first PWA
%   - Deploy: Docker Compose, Nginx host-level, Let's Encrypt HTTPS
%   - Reporting: ReportLab (PDF), per-task report templates
% Justify each choice briefly (why Daphne over uvicorn, why aiortc over
% mediasoup, etc.).

\section{Software Architecture}
\label{sec:impl-architecture}
% TODO. Suggested coverage:
%   - High-level diagram (web/db/redis containers + WebRTC ports + nginx host)
%   - Django app split (apps/sessions, apps/turns, apps/asr, apps/tts,
%     apps/moderation, apps/webrtc, apps/reports, apps/tasks)
%   - Session lifecycle: LOBBY -> INDIVIDUAL_RANKING -> ACTIVE -> CONCLUSION -> CLOSED
%   - Task plugin system (registry, base.py, four implementations)
%   - WebSocket endpoints: /ws/sessions/, /ws/turns/, /ws/webrtc/
%   - Turn-taking state machine (IDLE / HUMAN_SPEAKING / AI_SPEAKING /
%     AI_INTRODUCING) and reservation window
%   - Audio forwarding hub (no mixing, per-peer outbound queue,
%     backpressure via drop chunk)
%   - ASR pipeline (gating, OpenAI Realtime WebSocket, transcript persistence)
%   - TTS pipeline (gpt-4o-mini-tts, resample 24k->48k)
%   - Moderation orchestrator (triggers, LLM, intervention decisions)
%   - No-moderator mode (control arm) — guards in coordinator only
%   - GDPR consent flow, GDPR-compliant data retention

\section{UML Flow Diagrams}
\label{sec:impl-uml}
% TODO. Suggested diagrams:
%   - Sequence diagram: human turn lifecycle (frontend -> turns WS -> ASR ->
%     moderation orchestrator -> optional TTS -> back to IDLE)
%   - State diagram: TurnManager states + reservation window
%   - Sequence diagram: AI moderator intervention triggered by NO_PUSH timer
%   - Sequence diagram: forced conclusion at TIMER_30 expiry
%   - Activity diagram: session lifecycle with task-specific branches
%     (individual ranking phase only for survival tasks)

\chapter{Pilot Testing \& Design Changes}
\label{ch:pilot}

% Structure mirrors the host lab's thesis format (3 sections):
%   5.1 Study
%   5.2 Results and Discussion
%   5.3 Design Changes

This chapter reports the empirical evaluation of AIutami via pilot
studies and the design changes triggered by their findings.
Section~\ref{sec:pilot-study} describes the experimental study design.
Section~\ref{sec:pilot-results} presents quantitative and qualitative
results and discusses their implications.
Section~\ref{sec:pilot-design-changes} documents the changes made to
the system in response to the pilot findings, in line with the
iterative co-design methodology adopted by the host lab.

\section{Study}
\label{sec:pilot-study}
% TODO. Suggested coverage:
%   - Tasks: NASA Moon Survival (Hall, 1962; Hall & Watson, 1970) and
%     Lost at Sea (Nemiroff & Pasmore, 1975). 15-item ranking with
%     expert reference. Rationale for using both (within-subject design).
%   - Experimental design: within-subjects, counterbalanced
%     (moderator ON vs OFF; task order ON-first vs OFF-first).
%   - Participants: target N, recruitment, demographics, exclusion criteria,
%     network/device requirements (lessons from the 2026-05-11 pilot).
%   - Procedure: pre-questionnaire, individual ranking phase, group
%     discussion, post-questionnaire (SUS, custom UX scale).
%   - Metrics:
%       * Error score (sum of |team_rank - expert_rank|, 0-112)
%       * Synergy score (team vs best individual; team vs average individual)
%       * Gini index on speaking_time_per_participant
%       * SUS (System Usability Scale)
%       * Custom UX items (perceived helpfulness/intrusiveness of the moderator)

\section{Results and Discussion}
\label{sec:pilot-results}
% TODO. Suggested coverage:
%   - Quantitative results per metric (error score, synergy, Gini, SUS).
%     Statistical analysis (paired tests for within-subject comparisons;
%     check normality, choose parametric/non-parametric).
%   - Qualitative findings (open answers from post-questionnaire, observed
%     interaction patterns).
%   - Discussion: does the AI moderator improve performance? participation
%     balance? user experience? In which direction and by how much?
%   - Threats to validity (sample size, demographics, network quality,
%     credit availability for ASR/TTS during pilots).

\section{Design Changes}
\label{sec:pilot-design-changes}
% TODO. List the changes applied to AIutami in response to pilot findings:
% bugs surfaced and fixed, requirements refined, UX/UI adjustments,
% updated participant instructions, etc.

\chapter{Conclusion \& Future Works}
\label{ch:conclusion}

% Structure mirrors the host lab's thesis format (2 sections):
%   6.1 Conclusion
%   6.2 Future Works

\section{Conclusion}
\label{sec:concl-summary}
% TODO. Recap of the work:
%   - Summary of contributions (architecture; moderation strategies;
%     empirical evaluation), echoing the goals stated in Chapter 1.
%   - Direct answers to the research questions.
%   - Brief reflection on the relevance/impact of the findings.
%   - Limitations of the present work (small sample, network constraints
%     during pilots, etc.).

\section{Future Works}
\label{sec:concl-future}
% TODO. Open directions suggested by the work:
%   - Larger empirical study with a representative sample
%   - Additional tasks beyond NASA Moon and Lost at Sea (other survival
%     decision tasks; non-decision multi-party tasks)
%   - Predictive turn-taking (Voice Activity Projection) to replace the
%     explicit reservation window with implicit floor-management
%   - Emotion-aware moderation (currently the moderator is affect-blind)
%   - Multilingual support (system is currently Italian-only)
%   - Per-task fine-tuning of the moderator (e.g. specialised prompts for
%     different group sizes or for tasks with no expert ranking)


%-------------------------------------------------------------------------
%	BIBLIOGRAPHY
%-------------------------------------------------------------------------

\addtocontents{toc}{\vspace{2em}}
\bibliography{Thesis_bibliography}

%-------------------------------------------------------------------------
%	APPENDICES
%-------------------------------------------------------------------------

\cleardoublepage
\addtocontents{toc}{\vspace{2em}}
\appendix

\chapter{Questionnaires}
\label{app:questionnaires}

% TODO: paste the pre- and post-session questionnaires used in the empirical
% evaluation (Italian and/or English). Include the SUS items, NASA-TLX
% (if used), demographics, and any custom UX items.

% \input{Chapters/appendix_b_raw_data}  % TODO: enable when data is collected

% LIST OF FIGURES
\listoffigures

% LIST OF TABLES
\listoftables

% ACKNOWLEDGEMENTS
\chapter*{Acknowledgements}
% TODO: write at the end.

\cleardoublepage

\end{document}
